% section: 数学的帰納法

  \section{数学的帰納法}
    自然数$n$についての主張$P(n)$が,すべての$n\ge1$で成り立つことを示したいとする.
    このとき,以下の2つを示せばよい.
    \begin{enumerate}[label=(\roman*),parsep=3pt]
      \item $P(1)$が成り立つ.
      \item ある$n$で$P(n)$が成り立つと仮定すると,\,$P(n+1)$も成り立つ.
    \end{enumerate}
    この2つが示されれば,\,(i)より$P(1)$が成り立ち,\,(ii)より$P(1)\Rightarrow P(2)$,\,$P(2)\Rightarrow P(3)$,\,$\ldots$と次々に成り立つことが分かるので,すべての$n\ge1$について$P(n)$が成り立つ.
    ドミノ倒しに例えられることが多い.\,(i)は最初のドミノを倒すこと,\,(ii)は「あるドミノが倒れれば次のドミノも倒れる」ように並べることに対応する.

    本書では,漸化式から予想した一般項の形を確かめる場面や,三角関数の加法定理から漸化式を導く場面(3.2や3.6.3など)で数学的帰納法を用いる.